\documentclass[11pt]{article}
\usepackage{amsmath,amssymb,amsthm}
\usepackage{algorithm,algorithmic}
\usepackage{enumitem}
\usepackage{geometry}
\geometry{margin=1in}
\newtheorem{theorem}{Theorem}
\newtheorem{lemma}{Lemma}
\newtheorem{assumption}{Assumption}
\newcommand{\E}{\mathbb{E}}
\newcommand{\R}{\mathbb{R}}
\newcommand{\norm}[1]{\left\|#1\right\|}
\begin{document}

\section*{Randomized Block Coordinate Descent (RBCD)}

\section*{Mathematical Formulation}
Let $f:\R^{d}\rightarrow\R$ be a convex differentiable function.
Assume a partition of the coordinate set $\{1,\dots ,d\}$ into $B$ (non‑empty) blocks 
$\mathcal{B}_{1},\dots ,\mathcal{B}_{B}$.
For a vector $w\in\R^{d}$ we denote by $w_{\mathcal{B}_{i}}$ the subvector
containing the components in block $\mathcal{B}_{i}$ and by 
$\nabla_{\mathcal{B}_{i}} f(w)$ the corresponding partial gradient.

At iteration $t$ a block index $i_{t}\in\{1,\dots ,B\}$ is drawn 
independently from a distribution $P$ (typically uniform:
$P(i)=1/B$).  
Given a stepsize $\eta_{t}>0$, the update reads
\[
w^{t+1}=w^{t}-\eta_{t}\,U_{i_{t}}\nabla_{\mathcal{B}_{i_{t}}} f\bigl(w^{t}\bigr),
\tag{1}
\]
where $U_{i}$ is the $d\times |\mathcal{B}_{i}|$ zero‑padding matrix that
embeds a block vector into $\R^{d}$ (i.e. $U_{i}v$ has the components of $v$ in
the positions of $\mathcal{B}_{i}$ and zeros elsewhere).

\section*{Pseudocode}
\begin{algorithm}[H]
\caption{Randomized Block Coordinate Descent}
\begin{algorithmic}[1]
\STATE \textbf{Input:} initial point $w^{0}\in\R^{d}$, stepsize schedule
$\{\eta_{t}\}_{t\ge0}$, block partition $\{\mathcal{B}_{i}\}_{i=1}^{B}$.
\FOR{$t=0,1,2,\dots$}
    \STATE Sample block $i_{t}\sim P$ (uniform $\Rightarrow P(i)=1/B$).
    \STATE Compute partial gradient $g^{t}= \nabla_{\mathcal{B}_{i_{t}}} f(w^{t})$.
    \STATE $w^{t+1}= w^{t}-\eta_{t}\,U_{i_{t}} g^{t}$.
\ENDFOR
\STATE \textbf{Output:} last iterate $w^{T}$ or average $\bar w^{T}=\frac1T\sum_{t=0}^{T-1} w^{t}$.
\end{algorithmic}
\end{algorithm}

\section*{Dry Run Example}
Consider the univariate function $f(w)=(w-3)^{2}$.
Here $d=1$, $B=1$, the only block is $\mathcal{B}_{1}=\{1\}$.
Take stepsize $\eta=0.1$ and initialise $w^{0}=0$.

\begin{tabular}{c|c|c|c}
$t$ & $g^{t}= \nabla f(w^{t})$ & $w^{t+1}=w^{t}-\eta g^{t}$ & $f(w^{t+1})$\\ \hline
0 & $2(w^{0}-3) = -6$ & $0-0.1(-6)=0.6$ & $(0.6-3)^{2}=5.76$\\
1 & $2(0.6-3)=-4.8$ & $0.6-0.1(-4.8)=1.08$ & $(1.08-3)^{2}=3.6864$\\
2 & $2(1.08-3)=-3.84$ & $1.08-0.1(-3.84)=1.464$ & $(1.464-3)^{2}=2.3620$\\
\end{tabular}

The objective value decreases at each iteration, illustrating the
behaviour of (1) in the simplest setting.

\newpage

\section*{Assumptions}
\begin{assumption}[Block‑wise Lipschitz smoothness]\label{ass:smooth}
For each block $i\in\{1,\dots ,B\}$ there exists $L_{i}>0$ such that for all
$w,\tilde w\in\R^{d}$ differing only on block $i$,
\[
f(\tilde w)\le f(w)+\langle \nabla_{\mathcal{B}_{i}} f(w),\tilde w_{\mathcal{B}_{i}}-w_{\mathcal{B}_{i}}\rangle
+\frac{L_{i}}{2}\norm{\tilde w_{\mathcal{B}_{i}}-w_{\mathcal{B}_{i}}}^{2}.
\tag{2}
\]
\end{assumption}
\begin{assumption}[Convexity]\label{ass:convex}
$f$ is convex, i.e. for all $w,\tilde w\in\R^{d}$
\[
f(\tilde w)\ge f(w)+\langle \nabla f(w),\tilde w-w\rangle .
\tag{3}
\]
\end{assumption}
\begin{assumption}[Finite optimum]\label{ass:opt}
The set of minimisers $\mathcal{W}^{*}=\arg\min_{w}f(w)$ is non‑empty and we denote a
fixed optimal point by $w^{*}\in\mathcal{W}^{*}$.
\end{assumption}
\begin{assumption}[Stepsize]\label{ass:steps}
The stepsize is constant and chosen as
\[
\eta_{t}\equiv \eta=\frac{1}{L_{\max}},\qquad 
L_{\max}:=\max_{i}L_{i}.
\tag{4}
\]
\end{assumption}

\section*{Main Convergence Theorem}
\begin{theorem}[Expected $O(1/t)$ convergence]\label{thm:rate}
Under Assumptions \ref{ass:smooth}–\ref{ass:steps},
let $\{w^{t}\}_{t\ge0}$ be generated by the RBCD algorithm.
Define the averaged iterate
$\displaystyle \bar w^{T}:=\frac1T\sum_{t=0}^{T-1}w^{t}$.
Then for every $T\ge1$,
\[
\E\!\left[f(\bar w^{T})-f(w^{*})\right]
\;\le\;
\frac{2\,B\,L_{\max}\, \norm{w^{0}-w^{*}}^{2}}{T}.
\tag{5}
\]
Consequently $\E[f(w^{t})-f(w^{*})]=O(1/t)$.
\end{theorem}

\section*{Theoretical Properties}
\begin{itemize}[leftmargin=*]
\item \textbf{Stability.} The algorithm is a contraction in expectation
because the stepsize $\eta=1/L_{\max}$ guarantees that the term
$\frac{L_{i}\eta^{2}}{2}$ in the smoothness inequality (2) does not exceed
$\frac{\eta}{2}$.
\item \textbf{Hyperparameter sensitivity.}
The bound (5) depends linearly on $L_{\max}$ and on the number of blocks $B$.
Choosing a smaller stepsize $\eta<1/L_{\max}$ yields a larger constant
but does not affect the $O(1/T)$ rate; a larger stepsize would break the
inequality in Lemma~\ref{lem:descent}.
\item \textbf{Comparison to full‑gradient GD.}
If $B=1$ (i.e. a single block), (5) reduces to the classical bound
$\frac{2L\norm{w^{0}-w^{*}}^{2}}{T}$ for gradient descent.
When $B\gg1$ the factor $B$ appears, reflecting the fact that only a
$1/B$ fraction of coordinates is updated per iteration.
\end{itemize}

\section*{Discussion}
The theorem shows that randomised block updates retain the optimal
sublinear rate for convex smooth optimisation while potentially reducing
per‑iteration cost by a factor proportional to the block size. Extensions
to non‑uniform sampling, accelerated schemes, or strongly convex functions
follow by analogous modifications of the proof.

\newpage

\section*{Preliminaries (Definitions and Lemmas)}
\begin{lemma}[One‑step expected descent]\label{lem:descent}
Let $w^{t}$ be the current iterate and $i_{t}$ the sampled block.
With stepsize $\eta=1/L_{\max}$,
\[
\E_{i_{t}}\!\bigl[\,f(w^{t+1})\mid w^{t}\bigr]
\le f(w^{t})-\frac{\eta}{2B}\,\norm{\nabla f(w^{t})}^{2}.
\tag{6}
\]
\end{lemma}
\begin{proof}
Fix $w\equiv w^{t}$ and denote $g_{i}:=\nabla_{\mathcal{B}_{i}} f(w)$.
From block‑wise smoothness (2) with $\tilde w = w -\eta U_{i}g_{i}$ we obtain
\[
f\bigl(w-\eta U_{i}g_{i}\bigr)
\le f(w)-\eta\langle g_{i},g_{i}\rangle
+\frac{L_{i}\eta^{2}}{2}\norm{g_{i}}^{2}.
\tag{7}
\]
Since $\eta=1/L_{\max}\le 1/L_{i}$ for every $i$, we have $L_{i}\eta^{2}\le\eta$,
hence the third term in (7) is bounded by $\frac{\eta}{2}\norm{g_{i}}^{2}$.
Thus,
\[
f\bigl(w-\eta U_{i}g_{i}\bigr)
\le f(w)-\frac{\eta}{2}\norm{g_{i}}^{2}.
\tag{8}
\]
Taking expectation over the uniformly drawn block $i_{t}$,
\[
\E_{i_{t}}\!\bigl[f(w^{t+1})\mid w^{t}\bigr]
\le f(w^{t})-\frac{\eta}{2B}\sum_{i=1}^{B}\norm{g_{i}}^{2}
= f(w^{t})-\frac{\eta}{2B}\norm{\nabla f(w^{t})}^{2}.
\tag{9}
\]
Equation (9) is exactly (6). \hfill\qed
\end{proof}

\section*{Main Proof (Step‑by‑Step Derivation)}
We prove Theorem~\ref{thm:rate}.  
Let $w^{*}$ be any optimal solution (Assumption~\ref{ass:opt}).  
Define the squared distance $D^{t}:=\norm{w^{t}-w^{*}}^{2}$.

\begin{enumerate}
\item[Step 1.] Expand $D^{t+1}$ using the update (1):
\[
\begin{aligned}
D^{t+1}
&=\norm{w^{t}-\eta U_{i_{t}}g_{i_{t}}-w^{*}}^{2}\\
&=\norm{w^{t}-w^{*}}^{2}
-2\eta\langle U_{i_{t}}g_{i_{t}},\,w^{t}-w^{*}\rangle
+\eta^{2}\norm{U_{i_{t}}g_{i_{t}}}^{2}.
\end{aligned}
\tag{10}
\]
\item[Step 2.] Because $U_{i_{t}}$ only inserts zeros outside block $i_{t}$,
\[
\langle U_{i_{t}}g_{i_{t}},\,w^{t}-w^{*}\rangle
= \langle g_{i_{t}},\,w^{t}_{\mathcal{B}_{i_{t}}}-w^{*}_{\mathcal{B}_{i_{t}}}\rangle,
\qquad
\norm{U_{i_{t}}g_{i_{t}}}^{2}= \norm{g_{i_{t}}}^{2}.
\tag{11}
\]

\item[Step 3.] Apply convexity (3) with $\tilde w = w^{*}$:
\[
f(w^{t})-f(w^{*})\le\langle \nabla f(w^{t}),\,w^{t}-w^{*}\rangle.
\tag{12}
\]
Since $\nabla f(w^{t})$ concatenates all block gradients,
\[
\langle \nabla f(w^{t}),\,w^{t}-w^{*}\rangle
= \sum_{i=1}^{B}\langle g_{i},\,w^{t}_{\mathcal{B}_{i}}-w^{*}_{\mathcal{B}_{i}}\rangle .
\tag{13}
\]

\item[Step 4.] Take expectation of (10) conditioned on $w^{t}$ and use (11):
\[
\begin{aligned}
\E_{i_{t}}[D^{t+1}\mid w^{t}]
&= D^{t}
-2\eta\,\E_{i_{t}}\!\bigl[\langle g_{i_{t}},\,w^{t}_{\mathcal{B}_{i_{t}}}-w^{*}_{\mathcal{B}_{i_{t}}}\rangle\bigr]
+\eta^{2}\E_{i_{t}}\!\bigl[\norm{g_{i_{t}}}^{2}\bigr]\\
&= D^{t}
-\frac{2\eta}{B}\sum_{i=1}^{B}\langle g_{i},\,w^{t}_{\mathcal{B}_{i}}-w^{*}_{\mathcal{B}_{i}}\rangle
+\frac{\eta^{2}}{B}\sum_{i=1}^{B}\norm{g_{i}}^{2}.
\end{aligned}
\tag{14}
\]

\item[Step 5.] Substitute the sum from (13) and rewrite (14):
\[
\E_{i_{t}}[D^{t+1}\mid w^{t}]
= D^{t}
-\frac{2\eta}{B}\bigl(f(w^{t})-f(w^{*})\bigr)
+\frac{\eta^{2}}{B}\norm{\nabla f(w^{t})}^{2}.
\tag{15}
\]

\item[Step 6.] Use the stepsize choice $\eta=1/L_{\max}$ and the inequality
$\norm{\nabla f(w^{t})}^{2}\le 2L_{\max}\bigl(f(w^{t})-f(w^{*})\bigr)$,
which follows from smoothness of each block (see Lemma below).  
\begin{lemma}\label{lem:smooth-grad}
For any $w$,
\[
\norm{\nabla f(w)}^{2}\le 2L_{\max}\bigl(f(w)-f(w^{*})\bigr).
\]
\end{lemma}
\begin{proof}
From block smoothness (2) with $\tilde w = w^{*}$ (note $f(w^{*})\le f(w)$) we have
\[
f(w^{*})\ge f(w)+\langle \nabla f(w),\,w^{*}-w\rangle
-\frac{L_{\max}}{2}\norm{w-w^{*}}^{2}.
\]
Rearranging and applying Cauchy–Schwarz yields
\[
\langle \nabla f(w),\,w-w^{*}\rangle \ge f(w)-f(w^{*})-\frac{L_{\max}}{2}\norm{w-w^{*}}^{2}.
\]
By convexity, $\langle \nabla f(w),\,w-w^{*}\rangle\le \norm{\nabla f(w)}\norm{w-w^{*}}$.
Optimising over $\norm{w-w^{*}}$ gives the claimed bound. \hfill\qed
\end{proof}
Applying Lemma~\ref{lem:smooth-grad} to (15) we obtain
\[
\E_{i_{t}}[D^{t+1}\mid w^{t}]
\le D^{t}
-\frac{2\eta}{B}\bigl(f(w^{t})-f(w^{*})\bigr)
+\frac{2\eta^{2}L_{\max}}{B}\bigl(f(w^{t})-f(w^{*})\bigr).
\tag{16}
\]

\item[Step 7.] Since $\eta=1/L_{\max}$, the coefficient of the last term equals
$2\eta^{2}L_{\max}=2\eta/L_{\max}=2/L_{\max}^{2}\cdot L_{\max}=2\eta$,
hence (16) simplifies to
\[
\E_{i_{t}}[D^{t+1}\mid w^{t}]
\le D^{t}
-\frac{\eta}{B}\bigl(f(w^{t})-f(w^{*})\bigr).
\tag{17}
\]

\item[Step 8.] Take total expectation and sum (17) over $t=0,\dots ,T-1$:
\[
\begin{aligned}
\E[D^{T}]
&\le D^{0}
-\frac{\eta}{B}\sum_{t=0}^{T-1}\E\bigl[f(w^{t})-f(w^{*})\bigr].
\end{aligned}
\tag{18}
\]

\item[Step 9.] Rearrange (18) to isolate the average optimality gap:
\[
\frac{1}{T}\sum_{t=0}^{T-1}\E\bigl[f(w^{t})-f(w^{*})\bigr]
\le \frac{B\,D^{0}}{\eta\,T}
= \frac{B\,L_{\max}\norm{w^{0}-w^{*}}^{2}}{T}.
\tag{19}
\]

\item[Step 10.] By convexity, the function value at the averaged iterate
$\bar w^{T}$ does not exceed the average of the function values:
\[
f(\bar w^{T})\le \frac{1}{T}\sum_{t=0}^{T-1}f(w^{t}).
\tag{20}
\]
Taking expectation and combining (19)–(20) yields exactly the bound (5):
\[
\E\!\left[f(\bar w^{T})-f(w^{*})\right]
\le \frac{2\,B\,L_{\max}\norm{w^{0}-w^{*}}^{2}}{T}.
\]
\hfill\qed
\end{enumerate}

\section*{Rate Derivation (With Constants)}
The constant in (5) can be expressed as
\[
C:=2\,B\,L_{\max}\norm{w^{0}-w^{*}}^{2},
\]
so the expected suboptimality decays as $C/T$.
If $f$ is additionally $\mu$‑strongly convex,
the same proof with the stronger inequality
$\norm{\nabla f(w)}^{2}\ge 2\mu\bigl(f(w)-f(w^{*})\bigr)$
produces a linear rate
$\E[f(w^{t})-f(w^{*})]\le (1-\frac{\mu}{B L_{\max}})^{t}
\bigl(f(w^{0})-f(w^{*})\bigr)$.

\section*{Special Cases}
\begin{itemize}[leftmargin=*]
\item \textbf{Uniform vs. non‑uniform sampling.}
If block $i$ is chosen with probability $p_{i}>0$, the factor $B$ in (5)
is replaced by $\sum_{i}p_{i}^{-1}$, recovering the importance‑sampling
bounds found in the literature.
\item \textbf{Full‑gradient GD.}
When $B=1$, the algorithm reduces to standard gradient descent and (5)
coincides with the classical $2L\norm{w^{0}-w^{*}}^{2}/T$ bound.
\item \textbf{Parallel block updates.}
If $k$ disjoint blocks are updated simultaneously, the effective
stepsize can be increased to $k\eta$, yielding a bound with $B/k$ in place
of $B$.
\end{itemize}

\end{document}